\clearpage
\section{Certified context occurrences and graph-context congruence}
\label{sec:copying-certified-contexts}
Every annotated lambda-context now has an actual sole-hole address in its
unchanged Figure~16 translation. Its binding-free certificate was proved in
Section~\ref{sec:copying-context-binding}. We also prove composition of
addressed graph contexts, associativity of filling up to recursive graph
isomorphism, and graph-context congruence for an execution that respects that
isomorphism. The latter is an explicit semantic hypothesis to be discharged
by the copying machine, not an axiom of this development.

These results do not yet prove the paper's Lemma~11.3. That requires a further
structural bridge (whose unary case is subsequently proved in
Section~\ref{sec:copying-expansion-plug}): filling this certified graph with a suitably padded term
translation must agree, up to recursive isomorphism, with translating the
filled term. The constructor-interpreter equality from
Section~\ref{sec:copying-observations} is not this bridge: it concerns
re-evaluation of graph constructors, rather than addressed replacement.
Appendix~B.1, item~3 therefore remains open.

\subsection{Exact occurrence inventory}
\termdef{copying-single-hole-inventory}


Fix a named hole $h$, including its ordered input and output interfaces.
A primitive label has the single-hole property precisely when it is that
hole. A box label has the property precisely when exactly one edge of its
body has the property recursively and every other body edge is globally
hole-free. A graph has the property when exactly one outer edge has the
label property and all its other edges are globally hole-free. This counts
occurrences, including occurrences below boxes; repeated copies of the same
name do not count as one occurrence.

\lean{Copying/SingleHole}{11}{16}{SingleHole.SingleHoleLabel}


\lean{Copying/SingleHole}{17}{21}{SingleHole.SingleHole}


Interface equality transports an existing addressed site, and equality
elimination preserves its sole-hole certificate. Induction on depth then
turns the occurrence inventory into an actual site with exactly $h$'s
interfaces. At a primitive occurrence, choose its shallow edge and transport
the interface equations obtained from its label. At a box, use the induction
hypothesis in its body and extend the address through the box. The other-edge
conditions are precisely those required by the recursive sole-hole predicate.

The property is unchanged by an outer interface cast. A hole corolla has it
at every allowance: its only edge is the hole, and any proposed different
edge contradicts the singleton edge enumeration.

\lean{Copying/SingleHole}{22}{24}{SingleHole.Site.castInterfaces}


\lean{Copying/SingleHole}{25}{32}{SingleHole.SoleHole.castInterfaces}


\lean{Copying/SingleHole}{33}{59}{SingleHole.site}


\lean{Copying/SingleHole}{60}{64}{SingleHole.castInterfaces_iff}


\lean{Copying/SingleHole}{65}{69}{SingleHole.corolla}


\subsection{Preservation by the original assembly operations}

An edge bijection transports the inventory provided it preserves and
reflects both marked labels and hole-free labels. For the forward direction,
map the marked edge and pull each other edge back through the inverse;
injectivity shows that the inverse is different from the original marked
edge. The reverse argument is symmetric.

For an ordered parallel family, choose the marked edge in its unique marked
component. Another edge either belongs to that component, where the existing
certificate applies, or to another component, which is hole-free. Finite
expansion has the same tagged edge inventory and uses the same argument.
Serial composition is expansion of a two-edge pattern, so the left and right
versions follow by choosing the appropriate Boolean tag. Enclosing the graph
in a box makes that box the unique marked outer edge; the recursive label
condition is exactly the original graph inventory.

\lean{Copying/SingleHoleAssembly}{10}{32}{SingleHoleAssembly.transport}


\lean{Copying/SingleHoleAssembly}{33}{46}{SingleHoleAssembly.parallel}


\lean{Copying/SingleHoleAssembly}{47}{59}{SingleHoleAssembly.expand}


\lean{Copying/SingleHoleAssembly}{60}{67}{SingleHoleAssembly.serial_left}


\lean{Copying/SingleHoleAssembly}{68}{75}{SingleHoleAssembly.serial_right}


\lean{Copying/SingleHoleAssembly}{76}{83}{SingleHoleAssembly.box}


Raising preserves and reflects hole-freedom of an individual edge label.
For a primitive, the label constructor is unchanged. For a box, the existing
recursive all-labels padding theorem supplies the equivalence for its body.
Induction on allowance proves the marked-label equivalence as well; the edge
bijection then transports the entire single-hole inventory. Depth equality
casts preserve it by equality elimination. Iterating the raising theorem
handles finite padding, and combining padding with the depth cast handles
padding to any specified larger allowance.

\lean{Copying/SingleHoleAssembly}{84}{99}{SingleHoleAssembly.freeLabel_raise}


\lean{Copying/SingleHoleAssembly}{100}{121}{SingleHoleAssembly.raise_iff}


\lean{Copying/SingleHoleAssembly}{122}{124}{SingleHoleAssembly.castDepth_iff}


\lean{Copying/SingleHoleAssembly}{125}{130}{SingleHoleAssembly.pad_iff}


\lean{Copying/SingleHoleAssembly}{131}{134}{SingleHoleAssembly.padTo_iff}


A translated ordered pair retains exactly the marked occurrence in either
component when the other is hole-free. The proof uses the parallel theorem
and the existing boundary cast; it leaves the middle-block exchange intact.
Operation corollas are hole-free. Consequently abstraction preserves the
inventory by casting and boxing the body, prefixing the lambda operation,
and appending the original hole-free single-use distributor. Unary operations
use the same prefix and suffix argument without boxing. Application and
binary operations prefix their operation edge to the certified pair and
append the existing two-use environment distributor. No distributor edge
is eliminated or identified in any of these proofs.

\lean{Copying/SingleHoleTranslation}{10}{19}{SingleHoleTranslation.pair_left}


\lean{Copying/SingleHoleTranslation}{20}{29}{SingleHoleTranslation.pair_right}


\lean{Copying/SingleHoleTranslation}{30}{32}{SingleHoleTranslation.operation_holeFree}


\lean{Copying/SingleHoleTranslation}{33}{41}{SingleHoleTranslation.abstraction}


\lean{Copying/SingleHoleTranslation}{42}{48}{SingleHoleTranslation.unary}


\lean{Copying/SingleHoleTranslation}{49}{55}{SingleHoleTranslation.application}


\lean{Copying/SingleHoleTranslation}{56}{64}{SingleHoleTranslation.binary}


Structural induction on the annotated term context proves preservation
for any supplied single-hole graph. The hole case is immediate; abstraction
and unary cases use their constructor lemmas. In application and binary
cases, the induction hypothesis supplies the marked component and the fixed
term translation is hole-free. Padding both to the prescribed maximum
allowance preserves these facts. Thus the constructor lemma applies to the
literal graph used by the translator.

Supplying a hole corolla proves the exact named-hole inventory of the context
translation. Extracting the site and combining its sole-hole certificate
with the previously proved binding-free theorem gives an existential
certificate at the annotated term and atom interface. Its local allowance
is existential because enclosing fixed subterms can force additional padding.
No occurrence-existence assumption remains.

\lean{Copying/SingleHoleTranslation}{65}{89}{SingleHoleTranslation.context_apply}


\lean{Copying/SingleHoleTranslation}{90}{95}{SingleHoleTranslation.context_graph}


\lean{Copying/SingleHoleTranslation}{96}{101}{SingleHoleTranslation.context_certificate}


\subsection{Composition of actual addressed contexts}
\termdef{copying-addressed-composition}


Replacing a fixed box by recursively isomorphic bodies gives recursively
isomorphic outer graphs. Keep the outer vertex and edge maps equal to the
identity; at the selected edge use the given body isomorphism, and at every
other edge use reflexivity. Induction on an addressed site then proves that
plugging recursively isomorphic graphs preserves recursive isomorphism:
use generic shallow replacement at the outer level and the box result below
an enclosing box.

Replacing the same box twice retains only the last body. This follows
pointwise from the conditional label replacement and does not change ports.
Changing an inserted box body commutes with shallow insertion up to a graph
isomorphism whose carrier maps are identities. Retained edges have the old
labels, inserted edges have the same changed labels, and both port maps
follow the same host or inserted-vertex cases.

\lean{Copying/DeepComposition}{13}{35}{DeepComposition.replaceBox_deepIso}


\lean{Copying/DeepComposition}{36}{43}{DeepComposition.plugAt_deepIso}


\lean{Copying/DeepComposition}{44}{56}{DeepComposition.replaceBox_overwrite}


\lean{Copying/DeepComposition}{57}{99}{DeepComposition.plug_replaceBox_iso}


Let $s$ be an address in $G$ and $t$ an address in a graph $H$ that fits $s$.
Their composite is an address in $G[s:=H]$. If $s$ is shallow, an outer
edge of $H$ receives the inserted-edge tag; a deeper $t$ retains its body
address under that tag. If $s$ descends through a box, recursively compose
inside its body and then extend through the updated box.

Filling the composite address with $A$ is recursively isomorphic to
$G[s:=H[t:=A]]$. When both addresses are shallow this is the generic finite
graph associativity isomorphism. A shallow outer address and an inner box
use the preceding box-insertion interchange. For an outer box, remove the
overwritten intermediate body and lift the induction hypothesis through
that box. This proves associativity for arbitrary finite nested addresses.

If both original sites are sole holes, so is the composite: retained edges
are hole-free by the outer certificate and inserted non-selected edges by
the inner certificate. Induction handles enclosing boxes. Finally, when
the outer hole has a nonempty input interface, the sourced-hole closure
theorem proves that the composite host is binding-free. The zero-input
qualification from Section~\ref{sec:copying-binding} is retained.

\lean{Copying/DeepComposition}{100}{113}{DeepComposition.Site.comp}


\lean{Copying/DeepComposition}{114}{133}{DeepComposition.plugAt_assoc}


\lean{Copying/DeepComposition}{134}{159}{DeepComposition.SoleHole.comp}


\lean{Copying/DeepComposition}{160}{167}{DeepComposition.context_comp_certified}


\subsection{The semantic congruence theorem and its hypothesis}
\termdef{copying-execution-isomorphism}


An open graph context packages a graph, an addressed term hole, its
sole-hole proof and its binding-free proof. Its outer environment may remain
open. Filling it with an open program uses actual addressed replacement;
focus-freedom follows from replacement closure and hole-freedom from the
sole-hole theorem.

Compose a closing context $C$ with an open graph context $D$ by inserting
$D$'s graph and using the composite site. The preceding proofs supply the
sole-hole condition and binding-freedom; the hole's singleton term input
supplies the required nonempty interface. Associativity gives a recursive
isomorphism between closing $G$ with this composite context and closing
$D[G]$ with $C$.

\lean{Copying/ContextCongruence}{11}{17}{ContextCongruence.GraphContext}


\lean{Copying/ContextCongruence}{18}{26}{ContextCongruence.fill}


\lean{Copying/ContextCongruence}{27}{37}{ContextCongruence.ClosingContext.compose}


\lean{Copying/ContextCongruence}{38}{47}{ContextCongruence.ClosingContext.compose_fill_iso}


The execution interface currently accepts an arbitrary initial-state map
on closed graph presentations. Such a map need not respect renaming of
internal carriers. We therefore state its representation obligation
explicitly: recursively isomorphic closed, focus-free, hole-free graphs
must have universally cost-related executions. By the earlier observation
theorem, this means that they agree on successful termination. The copying
machine of Appendix~B.1, item~2 must eventually prove this property. It is
a parameter of the present congruence theorem, not a newly postulated axiom
and not a proved property of a concrete copying execution.

Suppose $G$ refines $H$ under all closing graph contexts and execution has
this invariance property. To prove refinement of $D[G]$ by $D[H]$, choose
an arbitrary closing context $C$. Apply the hypothesis to the proved
closing context $C\circ D$. Execution invariance transports the source
and target along the two filling-associativity isomorphisms. Transitivity
then yields the required comparison under $C$. Applying this argument to
both refinements proves graph-context equivalence congruence. The theorem
uses the universal cost relation; no stronger quantitative congruence is
claimed here.

\lean{Copying/ContextCongruence}{48}{55}{ContextCongruence.RespectsIso}


\lean{Copying/ContextCongruence}{56}{65}{ContextCongruence.refines_context}


\lean{Copying/ContextCongruence}{66}{71}{ContextCongruence.equiv_context}


Choose the allowance and address supplied by the translated-context
certificate and package them as an open graph context. Its graph is
\emph{definitionally equal} to the existing translation. Classical choice
selects a certified address, not a replacement graph or a new translation.
The graph-context congruence theorem now applies directly to this package,
under the explicit execution invariance hypothesis. Its conclusion compares
actual fillings of that graph. It does not yet identify those fillings
with translations of filled terms.

\lean{Copying/CertifiedContexts}{10}{17}{CertifiedContexts.certified}


\lean{Copying/CertifiedContexts}{18}{23}{CertifiedContexts.certified_graph}


\lean{Copying/CertifiedContexts}{24}{30}{CertifiedContexts.certified_graph_congruence}


A negative regression places two occurrences of the same hole in a parallel
graph. Whichever occurrence is selected, the other violates the requirement
that all remaining labels be hole-free. Thus the certificate does not confuse
unique names with unique occurrences.

The nested reference context from the preceding section receives an actual
address, sole-hole proof and binding-free proof at every hole allowance and
name. A reflexivity proof checks that its certified graph is exactly the
old Figure~16 graph. These structural regressions do not assert any copying
machine execution behaviour.

\lean{Copying/CertifiedContextExamples}{10}{20}{CertifiedContextExamples.duplicate_holes_rejected}


\lean{Copying/CertifiedContextExamples}{21}{28}{CertifiedContextExamples.nestedReferenceContext_certified}


\lean{Copying/CertifiedContextExamples}{29}{32}{CertifiedContextExamples.nestedReferenceContext_graph_unchanged}


The full Lean build and all three axiom audits pass. The completed
sole-hole certification and graph-context composition obligations are removed
from item~3's to-do wording. Addressed filling correspondence, its depth-padding
compatibility, and the resulting term-context congruence remain to be proved;
the concrete execution must also discharge its representation laws. Neither
Proposition~12.2 nor item~3 is marked complete.
